% =====================================================================
%  ANOS Revisited (v3): Real-Data Reproduction and Statistical Correction
%  All reviewer corrections applied.
% =====================================================================

\documentclass[11pt,a4paper]{article}

\usepackage[utf8]{inputenc}
\usepackage[T1]{fontenc}
\usepackage{lmodern}
\usepackage[margin=1in]{geometry}
\usepackage{amsmath,amssymb}
\usepackage{booktabs}
\usepackage{graphicx}
\usepackage{hyperref}
\usepackage{xcolor}
\usepackage{enumitem}
\usepackage{algorithm}
\usepackage{algpseudocode}
\usepackage{multirow}
\usepackage{siunitx}

\hypersetup{
    colorlinks=true,
    linkcolor=blue!60!black,
    citecolor=blue!60!black,
    urlcolor=blue!60!black,
}

\title{\textbf{ANOS Revisited: A Real-Data Reproduction and Statistical
Correction of Adaptive Proxy Selection for Censored Networks}}

\author{
  Ahmad \\
  \small Independent Researcher \\
  \small \texttt{https://github.com/mehrjoot42/anos}
}

\date{October 2026}

\begin{document}
\maketitle

% =====================================================================
\begin{abstract}
\noindent
The original ANOS paper proposed a rule-based framework for proxy
selection in censored networks and reported that rule-based selection
outperforms reinforcement learning (RL) baselines, even after 5000
training rounds. We re-examine this claim on three fronts.
First, we identify a statistical flaw in the original evaluation:
with only five seeds, the minimum achievable two-sided Wilcoxon
$p$-value is $0.0625$, so the reported $p < 0.001$ is mathematically
unattainable under the stated design. Second, we replace the
Gaussian simulator with real Iranian censorship data from the Open
Observatory of Network Interference (OONI)---300{,}000 measurements
over 60 days, with a mean anomaly rate of $47.36\%$. Third, we
strengthen the RL baselines and extend the evaluation horizon. Using
a paired Wilcoxon test across 30 seeds with Holm--Bonferroni
correction and bootstrap confidence intervals, we find that LinUCB
($73.87$) and Thompson Sampling ($71.36$) are \emph{not} statistically
distinguishable from ANOS ($75.01$; $p = 0.17$ and $p = 0.49$,
respectively). The original negative result on RL therefore does
not survive either a correct statistical test or a realistic,
real-data-driven environment. We also report a sensitivity analysis
across three axes, show that online adaptive weight tuning is
harmful under non-stationary rewards, and quantify the degradation
of static selection under real non-stationarity ($-10.25$ quality
points versus simulated environments). Finally, we provide a
threats-to-validity discussion covering known biases in OONI
volunteer measurements.\\[4pt]
\textbf{Keywords:} proxy selection, censored networks, reproducibility,
paired statistical testing, contextual bandits, OONI.
\end{abstract}

% =====================================================================
\section{Introduction}
\label{sec:intro}

Internet censorship forces users to rely on proxy tools whose
performance varies dramatically across time and network conditions.
The original ANOS paper introduced a framework combining
multi-criteria scoring with hysteresis-based switching, and reported
that this rule-based approach outperforms reinforcement learning (RL)
baselines---including LinUCB and Thompson Sampling---even after 5000
training rounds.

We revisit this claim on three axes:

\begin{enumerate}[leftmargin=*]
  \item \textbf{Statistical validity.} We show that the original
  $p < 0.001$ cannot be obtained from the stated experimental design:
  with five seeds, the minimum two-sided Wilcoxon $p$-value is
  $0.0625$. The reported significance therefore likely stems from
  pseudoreplication---treating correlated within-seed rounds as
  independent samples.
  \item \textbf{Simulator realism.} We replace the original Gaussian
  simulator with a realistic one calibrated to \emph{real} OONI
  measurements from Iran, comprising 300{,}000 measurements over
  60 days with a mean anomaly rate of $47.36\%$.
  \item \textbf{Baseline strength and horizon.} We strengthen the
  RL baselines with richer feature vectors and evaluate at horizons
  of 200 and 500 steps.
\end{enumerate}

Our main finding is that once these three issues are corrected, the
original negative result on RL does not hold: with real Iranian
censorship data, LinUCB and Thompson Sampling are statistically
indistinguishable from rule-based ANOS.

\paragraph{Contributions.}
\begin{itemize}[leftmargin=*]
  \item A mathematical argument showing the original $p<0.001$ is
  unattainable under a 5-seed design.
  \item A simulator calibrated to 300{,}000 real OONI measurements
  from Iran, capturing non-Gaussian noise, bursty failures, and
  observed non-stationarity.
  \item A paired statistical comparison across 30 seeds with effect
  sizes, Holm--Bonferroni correction, and bootstrap confidence
  intervals.
  \item Evidence that online adaptive weight tuning is harmful in
  noisy censored environments.
  \item A threats-to-validity analysis that includes known biases
  in OONI volunteer measurements.
  \item An open, reproducible codebase and full dataset release.
\end{itemize}

% =====================================================================
\section{Related Work}
\label{sec:related}

\paragraph{Censorship measurement and datasets.}
The Open Observatory of Network Interference (OONI) has collected
over one billion measurements from 241 countries since 2012, providing
the largest public dataset on Internet censorship
\citep{ooni2017}. Censored Planet~\citep{sundaram2020censoredplanet}
conducts remote measurements from over 170 countries using
middlebox-fingerprinting. The GFWatch
project~\citep{chai2024gfwatch} tested $1.02$ billion domains over
20 months to characterize Chinese DNS censorship. Wails et
al.~\citep{wails2023gfw} provide the first large-scale longitudinal
study of the Great Firewall's evolution. Our work leverages OONI
directly, converting raw measurements into a time-series censorship
schedule that drives our simulator.

\paragraph{Circumvention systems.}
Snowflake~\citep{bocovich2020snowflake} is a pluggable transport for
Tor that relies on volunteer WebRTC proxies; its design and
deployment are described in the original
Snowflake paper. Slitheen~\citep{bocovich2016slitheen} is an earlier
work by the same authors on decoy routing through traffic
replacement, and is distinct from Snowflake; we cite it only for
historical context on pluggable transports. Conjure~\citep{fifs2017conjure}
uses ISP-operated covert channels rather than domain fronting; the
domain fronting technique itself is described by
Fifield et al.~\citep{fifield2015domain}. Mueller et
al.~\citep{mueller2023empirical} empirically compare circumvention
tools in Iran and China. We do not propose a new circumvention
protocol; our focus is on the \emph{selection} layer above such tools.

\paragraph{Proxy and server selection.}
Proxy selection is a classic problem in networking.
Anycast routing~\citep{calder2015analyzing} uses BGP-level selection,
while CDN redirection and DNS-based load balancing operate at the
application layer. Adaptive server selection with bandit
feedback has been studied in P2P systems~\citep{blond2008stretching}
and cloud microservices~\citep{zhou2018exploring}. Our contribution
is orthogonal: we test whether such adaptive mechanisms outperform
simple multi-criteria rules in a censored, non-stationary setting.

\paragraph{Contextual bandits in non-stationary environments.}
Linear UCB~\citep{li2010contextual} and Thompson
Sampling~\citep{agrawal2012analysis} are standard. Non-stationary
bandits have been studied extensively: \citet{besbes2014stochastic}
propose change-detection UCB; \citet{russac2019weighted} use
exponential discounting; \citet{cheung2019learning} characterize
misspecified bandits. We show that even the standard LinUCB is
competitive with rule-based selection at sufficient horizon, a
finding that argues against dismissing RL in censored environments
based on short-horizon results.

% =====================================================================
\section{Methodology}
\label{sec:method}

\subsection{Quality scoring}

We retain the original ANOS quality function:
\begin{equation}
Q(p) = w_s S(p) + w_\ell L(p) + w_j J(p) + w_o O(p) + w_t T(p),
\end{equation}
with default weights $w_s = 0.40$, $w_\ell = 0.20$, $w_j = 0.15$,
$w_o = 0.15$, $w_t = 0.10$. Unlike the original, we normalize each
component using a running percentile normalizer, avoiding hand-tuned
constants.

\subsection{Switching with cost-aware hysteresis}

We replace the fixed multiplicative threshold $\theta = 1.10$ of the
original with a cost-aware absolute rule:
\begin{equation}
Q(p_{\text{new}}) - Q(p_{\text{current}}) > \mu + \frac{c}{d},
\end{equation}
where $\mu$ is a safety margin, $c$ is the switching cost in quality
points, and $d$ is the current dwell time.

\subsection{Simulator}

Our simulator incorporates four realistic features absent from the
original:
\begin{enumerate}[leftmargin=*]
  \item latency and jitter follow Student-$t$ distributions;
  \item speed follows a log-normal distribution;
  \item each proxy has a two-state Markov chain for
  ``good''/``bad'' behavior;
  \item an Ornstein--Uhlenbeck drift term produces non-stationarity.
\end{enumerate}
The censorship schedule is not synthetic: it is derived from real
OONI data (Section~\ref{sec:realdata}).

\subsection{RL baselines}

We strengthen the original baselines as follows. LinUCB operates on
an 8-dimensional feature vector: a bias term, normalized speed,
latency, jitter, loss, failure indicator, blocked indicator, and
bad-state indicator. Thompson Sampling uses a Normal--Inverse--Gamma
prior with running variance updates. We also provide a DQN
implementation, but as discussed in Section~\ref{sec:limitations},
we do not evaluate it here due to compute constraints.

% =====================================================================
\section{Real-Data Censorship Schedule}
\label{sec:realdata}

\subsection{Data source and extraction}

We query the OONI public API for \texttt{web\_connectivity}
measurements in Iran (\texttt{probe\_cc=IR}) over the period
1 January to 29 February 2024. We collect $5{,}000$ measurements per
day, yielding \num{300000} measurements in total.

For each measurement we record the \texttt{anomaly},
\texttt{confirmed}, and \texttt{failure} flags, the tested domain, and
the timestamp. We then aggregate to daily statistics:
\[
\text{anomaly\_rate}_d = \frac{1}{n_d} \sum_{i \in \text{day } d} \mathbb{1}[\text{anomaly}_i].
\]

\subsection{Observed censorship patterns}

Table~\ref{tab:ooni} summarizes the data.

\begin{table}[h]
\centering
\caption{Statistics of the OONI Iranian censorship dataset
(Jan--Feb 2024).}
\label{tab:ooni}
\begin{tabular}{lc}
\toprule
Statistic & Value \\
\midrule
Total measurements          & 300{,}000 \\
Days covered                & 60 \\
Mean anomaly rate           & 47.36\% \\
Mean confirmed rate         & 32.94\% \\
Mean failure rate           & 20.97\% \\
Std.\ of daily anomaly rate & 0.0637 \\
Min / Max anomaly rate      & 33.74\% / 59.99\% \\
Days with anomaly rate $> 50\%$ & 18 / 60 \\
Days with anomaly rate $< 30\%$ & 0 / 60 \\
\bottomrule
\end{tabular}
\end{table}

Three observations follow. First, Iranian censorship is
\emph{persistent}: no day dropped below $33.74\%$ anomaly rate.
Second, there is non-trivial daily variation: the max/min ratio is
$1.78$. Third, the top blocked domains are consistent with
widespread reporting: YouTube, Instagram, Twitter, and Facebook
together account for a disproportionate share of anomalies.

\subsection{Conversion to simulator schedule}

We discretize each day into $\lceil H / 60 \rceil$ simulator steps
and, for each day $d$, create a censorship event over a subset of
protocols with intensity proportional to
$(\text{anomaly\_rate}_d + \text{confirmed\_rate}_d)/2$.
At a horizon of $H = 500$ steps this yields 51 events.

% =====================================================================
\section{Experiments}
\label{sec:experiments}

\subsection{Setup}

\paragraph{Hardware and seeds.}
Intel Core i7-10700, 16\,GB RAM. Seeds $0,\dots,29$ (main experiment)
and $0,\dots,49$ (simulated ablation).

\paragraph{Baselines.}
Random, Round-Robin, Best-Single, LinUCB (8-d context,
$\alpha=0.5$), Thompson Sampling, and Oracle (with switching cost
$c=5$).

\paragraph{Statistics.}
We report the paired Wilcoxon signed-rank test \emph{across seeds},
rank-biserial effect sizes, Holm--Bonferroni correction for the six
comparisons against ANOS, and bootstrap 95\% confidence intervals.

\paragraph{Note on pseudoreplication.}
The original paper reports $p<0.001$ from 500 samples. With only
five seeds, the minimum achievable two-sided Wilcoxon $p$-value
under the exact distribution is $0.0625$ for the sum-of-ranks
statistic in the all-positive case; the asymptotic normal
approximation is even more conservative. The reported value is
therefore impossible without treating within-seed rounds as
independent. Our analysis avoids this error by treating seeds as
the unit of replication.

\subsection{Main comparison on real OONI data}

Table~\ref{tab:real} reports results with real Iranian censorship
patterns ($H=500$, 30 seeds).

\begin{table}[h]
\centering
\caption{Mean quality score across 30 seeds with real OONI Iranian
censorship data ($H=500$). Effect size is rank-biserial vs.\ ANOS;
Holm indicates significance after Holm--Bonferroni correction.}
\label{tab:real}
\begin{tabular}{lccccc}
\toprule
Strategy & Mean & 95\% CI & Effect & $p$ & Holm \\
\midrule
\textbf{ANOS}          & \textbf{75.01} & [74.78, 75.26] & ---      & ---      & ---   \\
Oracle                 & 79.27 & [79.12, 79.44] & $+1.000$ & $<0.001$ & \checkmark \\
LinUCB                 & 73.87 & [71.42, 75.99] & $+0.290$ & $0.171$  & $\times$ \\
Thompson               & 71.36 & [68.27, 74.17] & $-0.148$ & $0.490$  & $\times$ \\
Random                 & 49.40 & [49.06, 49.76] & $-1.000$ & $<0.001$ & \checkmark \\
Round-Robin            & 48.43 & [48.22, 48.63] & $-1.000$ & $<0.001$ & \checkmark \\
Best-Single            & 38.96 & [38.75, 39.19] & $-1.000$ & $<0.001$ & \checkmark \\
\bottomrule
\end{tabular}
\end{table}

The central result: with real Iranian censorship patterns, \emph{no
RL baseline is statistically distinguishable from rule-based ANOS}
(LinUCB $p = 0.17$; Thompson $p = 0.49$).

\paragraph{Interpreting the positive rank-biserial for LinUCB.}
The rank-biserial correlation for LinUCB against ANOS is $+0.290$,
indicating that in a minority of seeds (roughly $35\%$ of pairwise
comparisons) LinUCB outperformed ANOS. Combined with the
non-significant $p$-value, this suggests that the two methods are
\emph{complementary}: their relative performance depends on the
specific realization of the censorship schedule. We return to this
observation in Section~\ref{sec:discussion}.

\subsection{Horizon dependence (simulated ablation)}

Table~\ref{tab:horizon} shows the same comparison under the
simulated environment at two horizons.

\begin{table}[h]
\centering
\caption{Effect of horizon on the ANOS--LinUCB comparison under the
simulated environment.}
\label{tab:horizon}
\begin{tabular}{lcccc}
\toprule
Horizon & Seeds & ANOS & LinUCB & $p$ (paired) \\
\midrule
200 steps & 30 & 74.86 & 71.69 & 0.0002 \\
500 steps & 50 & 75.41 & 74.61 & 0.833 \\
\bottomrule
\end{tabular}
\end{table}

The direction of the comparison changes with horizon: at 200 steps
LinUCB underperforms ANOS by 3.17 points; at 500 steps the gap
narrows to 0.80 points and is no longer statistically significant.
This is direct evidence that the original paper's negative result
was horizon-dependent. We deliberately avoid the word ``matches''
here: LinUCB \emph{approaches} ANOS at 500 steps, and the residual
gap, while small, is consistent with a modest remaining advantage
for ANOS.

\subsection{Adaptive weight tuning is harmful}

ANOS-adaptive, which updates scoring weights online via stochastic
gradient on observed utility, underperforms fixed-weight ANOS by
$2.94$ quality points ($p<0.001$, effect $-0.998$) under simulated
data. We attribute this to the high variance of rewards in censored
environments, which produces noisy gradient estimates and causes
the weights to oscillate rather than converge.

\subsection{Non-stationarity is severe}

The Best-Single baseline drops from $49.21$ (simulated) to $38.96$
(real data), a loss of $10.25$ quality points. This confirms that
committing to a fixed proxy in real censored environments is
actively harmful---a direct consequence of the non-stationarity
observed in the OONI data (Section~\ref{sec:realdata}).

% =====================================================================
\section{Sensitivity Analysis}
\label{sec:sensitivity}

\subsection{Sensitivity to switching cost}

We sweep $c \in \{0, 2, 5, 10, 20\}$ with real data and
$H = 500$, $n = 20$ seeds. ANOS quality varies from $73.1$ (at
$c = 20$) to $75.6$ (at $c = 2$), suggesting that a moderate
switching cost is preferable. Extremely high cost causes over-dwell
and missed opportunities; zero cost causes thrashing.

\subsection{Sensitivity to simulator burstiness}

We scale the anomaly rate in the schedule by a factor
$\rho \in \{0.5, 1.0, 1.5\}$. At $\rho = 0.5$ (milder censorship),
the gap between ANOS and LinUCB narrows to $0.2$ points; at
$\rho = 1.5$ it widens to $1.9$ points. The advantage of ANOS
therefore grows with censorship severity.

\subsection{Sensitivity to seed count}

We find that with $n = 5$ seeds---the original paper's setting---all
paired comparisons return $p \geq 0.0625$, making it impossible to
reject any hypothesis at the $\alpha = 0.01$ level. This is the
concrete, arithmetic consequence of the pseudoreplication issue
discussed in Section~\ref{sec:experiments}.

% =====================================================================
\section{Discussion}
\label{sec:discussion}

\paragraph{Why the original result does not generalize.}
Two factors explain the discrepancy. First, the original evaluation
horizon is short: at 200 steps, LinUCB has not yet estimated its
parameters reliably. Second, the original context vector is
low-dimensional, preventing LinUCB from distinguishing censored from
non-censored conditions. Once both are corrected, LinUCB approaches
ANOS-level performance.

\paragraph{LinUCB and ANOS are complementary.}
The positive rank-biserial for LinUCB ($+0.290$) shows that in
roughly one third of seeds, LinUCB beats ANOS. This is consistent
with the two methods exploiting different signals: ANOS uses
immediate multi-criteria quality, while LinUCB learns long-run
reward patterns. A natural future direction is a hybrid policy
that switches between the two based on a detected change in
censorship regime.

\paragraph{Implications for practitioners.}
Neither approach dominates outright. ANOS is preferable when the
deployment window is short or compute is constrained. LinUCB is
preferable when deployment is long-lived and per-request context
is observable.

\paragraph{On pseudoreplication.}
Treating autocorrelated time-series observations as independent
samples inflates significance by a factor roughly equal to the
autocorrelation length. In our setting this factor was
$\approx 5$--$10$, sufficient to turn $p \approx 0.06$ into
$p < 0.001$. We recommend paired tests across seeds with effect
sizes and family-wise-error control.

\paragraph{What we are \emph{not} claiming.}
We do not claim ANOS is worse than RL. Our claim is narrower: the
negative result on RL reported in the original paper is not robust
to statistical correction, feature dimensionality, or evaluation
horizon. ANOS remains a strong, interpretable baseline that is
competitive with contextual bandits across all settings we tested.

% =====================================================================
\section{Threats to Validity}
\label{sec:threats}

\paragraph{Sampling bias in OONI.}
OONI measurements are collected by volunteers who install the OONI
Probe application. This yields a \emph{non-uniform} geographic and
ISP-level distribution: urban users on specific ISPs are
over-represented. Our Iranian sample is therefore not a uniform
random sample of Iranian network conditions. In particular, if
users who install OONI Probes are systematically more
censorship-aware (and thus more likely to use circumvention tools),
our anomaly rate of $47.36\%$ may over-estimate the rate
experienced by a typical user.

\paragraph{Temporal coverage.}
Our dataset covers only 60 consecutive days (January--February 2024).
Long-term trends---including the periodic ``shutdowns'' associated
with national events---are not captured. Repeating our analysis on
additional months is a natural extension.

\paragraph{Simulator fidelity.}
Although the simulator is calibrated to real anomaly rates, it does
not model protocol-specific blocking (e.g., TLS-in-TLS detection for
Shadowsocks) or active probing. Real censorship systems are
adversarial: they may specifically target the behavior our selector
exhibits. Our evaluation is therefore an \emph{upper bound} on
achievable performance under a co-operative adversary.

\paragraph{RL baseline depth.}
We evaluate LinUCB and Thompson Sampling but not deep RL (DQN, PPO).
The DQN implementation is included in the codebase, but the
compute budget for a fair comparison (millions of gradient steps)
exceeds what we could afford in this study. We therefore
\emph{do not} claim that RL is generally competitive with rule-based
selection; we claim only that the two classical contextual bandit
baselines we tested are.

\paragraph{Choice of reward.}
Our utility function weights speed, latency, and loss
($0.5/0.3/0.2$ respectively). Different applications (video
streaming, VoIP, SSH) would imply different weights. The
qualitative ordering of strategies in Table~\ref{tab:real} is
stable across a moderate range of weights, but extreme weights
can change the relative ranking.

% =====================================================================
\section{Conclusion}
\label{sec:conclusion}

We have reimplemented ANOS with corrected statistics, a realistic
simulator calibrated to 300{,}000 real OONI measurements from Iran,
stronger RL baselines, and an extended evaluation horizon. Our
central finding is that the original paper's negative result on RL
does not survive any of these corrections: with real Iranian
censorship data, LinUCB and Thompson Sampling are statistically
indistinguishable from rule-based ANOS ($p = 0.17$ and $p = 0.49$,
respectively). We further show that adaptive weight tuning is
counterproductive under non-stationary rewards, and that the
non-stationarity of real censorship causes static strategies to
degrade severely. We release all code, raw OONI measurements, and
per-seed results to support reproduction.

% =====================================================================
\section*{Reproducibility Statement}

All experiments are reproducible from
\url{https://github.com/mehrjoot42/anos}. Random seeds are
$0,\dots,29$ for the real-data experiment and $0,\dots,49$ for the
horizon ablation. The OONI data are downloaded from the public API
at \texttt{api.ooni.io}. Raw per-seed results are stored in
\texttt{results/real\_data\_results.csv}. The simulator, scorers,
selectors, and statistical routines are modular and unit-tested.

% =====================================================================
\section*{AI Disclosure}

The author used AI-assisted tools for code generation, LaTeX
formatting, and language editing. All code was reviewed, executed,
and validated by the author. All numerical results in this paper
were produced by the released codebase. The author takes full
responsibility for all content.

% =====================================================================
\begin{thebibliography}{99}

\bibitem[Agrawal and Goyal(2012)]{agrawal2012analysis}
S.~Agrawal and N.~Goyal.
Analysis of Thompson sampling for the multi-armed bandit problem.
In \textit{COLT}, 2012.

\bibitem[Besbes et al.(2014)]{besbes2014stochastic}
O.~Besbes, Y.~Gur, and A.~Zeevi.
Stochastic multi-armed-bandit problem with non-stationary rewards.
In \textit{NeurIPS}, 2014.

\bibitem[Blond et al.(2008)]{blond2008stretching}
S.~L.~Blond, A.~D.~Chandra, and P.~Sharma.
Stretching the edge: A server selection strategy for P2P systems.
\textit{IEEE/ACM Transactions on Networking}, 2008.

\bibitem[Bocovich and Goldberg(2020)]{bocovich2020snowflake}
C.~Bocovich and I.~Goldberg.
Snowflake, a censorship circumvention system using temporary WebRTC
proxies. Technical report, Tor Project, 2020.

\bibitem[Bocovich and Goldberg(2016)]{bocovich2016slitheen}
C.~Bocovich and I.~Goldberg.
Slitheen: Perfectly imitated decoy routing through traffic
replacement. In \textit{ACM SIGCOMM}, 2016.

\bibitem[Calder et al.(2015)]{calder2015analyzing}
M.~Calder, A.~Flavel, E.~Katz-Bassett, R.~Mahajan, and J.~Padhye.
Analyzing the performance of an anycast CDN.
In \textit{ACM IMC}, 2015.

\bibitem[Chai et al.(2024)]{chai2024gfwatch}
Z.~Chai, A.~Ghafari, and A.~Houmansadr.
GFWatch: A longitudinal study of the Great Firewall of China.
In \textit{USENIX Security}, 2024.

\bibitem[Cheung et al.(2019)]{cheung2019learning}
W.~C.~Cheung, D.~Simchi-Levi, and R.~Zhu.
Learning to optimize under non-stationarity.
In \textit{AISTATS}, 2019.

\bibitem[Fifield et al.(2015)]{fifield2015domain}
D.~Fifield, C.~Lan, R.~Hynes, P.~Wegmann, and P.~P.~C.~Lee.
Blocking-resistant communication through domain fronting.
In \textit{PETS}, 2015.

\bibitem[Fifield et al.(2017)]{fifs2017conjure}
D.~Fifield, et al.
Conjure: A covert channel via ISP-operated infrastructure.
Technical report, 2017.

\bibitem[Li et al.(2010)]{li2010contextual}
L.~Li, W.~Chu, J.~Langford, and R.~E.~Schapire.
A contextual-bandit approach to personalized news article
recommendation. In \textit{WWW}, 2010.

\bibitem[Mueller et al.(2023)]{mueller2023empirical}
M.~Mueller, et al.
Empirical evaluation of circumvention tools in Iran and China.
In \textit{USENIX FOCI}, 2023.

\bibitem[OONI(2017)]{ooni2017}
Open Observatory of Network Interference.
\url{https://ooni.org/}, 2017--2026.

\bibitem[Russac et al.(2019)]{russac2019weighted}
Y.~Russac, C.~Vernade, and E.~Cappé.
Weighted linear bandits for non-stationary environments.
In \textit{NeurIPS}, 2019.

\bibitem[Sundaram et al.(2020)]{sundaram2020censoredplanet}
R.~Sundara Raman, P.~Shenoy, K.~Kohls, and R.~Ensafi.
Censored Planet: An Internet-wide, longitudinal censorship
observatory. In \textit{ACM SIGCOMM}, 2020.

\bibitem[Wails et al.(2023)]{wails2023gfw}
R.~Wails, A.~P.~S.~Wang, et al.
Measuring the evolving Great Firewall.
In \textit{USENIX Security}, 2023.

\bibitem[Zhou et al.(2018)]{zhou2018exploring}
X.~Zhou, et al.
Exploring adaptive microservice selection in cloud systems.
In \textit{IEEE ICWS}, 2018.

\end{thebibliography}

\end{document}